% =====================================================================
% PLANTILLA: Guía del estudiante (trimestral) -- Área de Matemáticas
% María Mercedes Cantillo · Colegio San Alberto Magno
% ---------------------------------------------------------------------
% Una guía por periodo (trimestre). Soporta los temas de las clases del
% periodo: cada clase (clase.tex) remite a los temas de esta guía.
%
% Estructura obligatoria:
%   1. \frase        frase introductoria (cita real, con fuente)
%   2. Introducción   objetivo del periodo + entrada al hilo conductor
%   3. DBA            enunciado oficial tomado de dba/ (nunca de memoria)
%   4. Marco teórico  soporte académico; aquí van las citas de científicos
%   5. \tema{...}     un bloque por tema, cada uno con:
%                      hilo -> explicación -> aplicacion -> ejercicios -> resumen
%   6. Cierre del hilo conductor
%   7. Evaluación     secciones institucionales (matriz, autoevaluación, seguimiento)
%   8. Referencias    toda cita del documento debe aparecer aquí
%
% Hilo conductor: una historia que atraviesa toda la guía (historia de la
% humanidad, o literatura: Harry Potter, Las crónicas de Narnia, El señor
% de los anillos...). Este ejemplo usa la historia del cálculo.
%
% Compilar:  latexmk -pdf guia-periodo-I-derivadas.tex
% =====================================================================
\documentclass[11pt]{article}
\makeatletter\def\input@path{{../estilo/}}\makeatother
\usepackage{mmcantillo}

% ======================== DATOS DE LA GUÍA ============================
\tipodocumento{Guía del estudiante}
\asignatura{Cálculo}
\titulo{La matemática del cambio: la derivada}
\grado{11°}
\periodo{I}
\hiloconductor{La carrera por el cálculo: de Galileo a la Luna}
% ======================================================================
% DBA: matematicas grado 11 · DBA 5 y DBA 8 (dba/matematicas/grados/grado11.tex)

\begin{document}

\encabezado

% ---------------------------------------------------------------------
% 1. Frase introductoria
% ---------------------------------------------------------------------
% verificado: https://portalegalileo.museogalileo.it/egjr.asp?c=36261 — Il Saggiatore (Roma, 1623), pasaje del «grandissimo libro» escrito «in lingua matematica»
\frase{La filosofía está escrita en ese grandísimo libro que tenemos
abierto ante los ojos, quiero decir, el universo, pero no se puede
entender si antes no se aprende a entender la lengua [\ldots] en que está
escrito. Está escrito en lengua matemática.}
{Galileo Galilei, \emph{El ensayador} (1623)}

% ---------------------------------------------------------------------
% 2. Introducción
% ---------------------------------------------------------------------
\seccion{Introducción}

Todo cambia: la posición de un dron, la temperatura de una taza de café,
el precio de un producto, la población de una ciudad. En este periodo
aprenderemos la herramienta que la humanidad inventó para \emph{medir el
cambio en un instante}: la \textbf{derivada}.

Nuestro hilo conductor es una historia real. A finales del siglo XVII, dos
hombres que nunca trabajaron juntos ---Isaac Newton en Inglaterra y
Gottfried Leibniz en Alemania--- inventaron el cálculo casi al mismo
tiempo, y terminaron enfrentados en una de las disputas más famosas de la
ciencia. Seguiremos a sus herederos: una marquesa francesa que tradujo a
Newton, una matemática que Einstein admiraba y una mujer que calculó a
mano la ruta de un astronauta. En cada tema daremos un paso de esa carrera.

\textbf{Al terminar este periodo podrás:}
\begin{itemize}
  \item interpretar la derivada como razón de cambio y como pendiente de
        la recta tangente;
  \item calcular derivadas de funciones polinómicas con las reglas básicas;
  \item usar la derivada para describir el movimiento y resolver problemas.
\end{itemize}

Este trimestre trabaja dos Derechos Básicos de Aprendizaje del Ministerio de
Educación Nacional~\cite{mendba}:

\begin{dba}[Grado 11 · DBA 5]
Interpreta la noción de derivada como razón de cambio y como valor de la
pendiente de la tangente a una curva y desarrolla métodos para hallar las
derivadas de algunas funciones básicas en contextos matemáticos y no
matemáticos.
\end{dba}

\begin{dba}[Grado 11 · DBA 8]
Encuentra derivadas de funciones, reconoce sus propiedades y las utiliza
para resolver problemas.
\end{dba}

% ---------------------------------------------------------------------
% 3. Marco teórico (las citas van aquí, con su referencia)
% ---------------------------------------------------------------------
\seccion{Marco teórico}

\subsubsection*{Un problema antiguo: medir lo que cambia}

Desde la Antigüedad se sabía calcular una velocidad \emph{promedio}:
distancia recorrida dividida entre el tiempo empleado. Lo difícil era la
velocidad \emph{en un instante}, cuando el tiempo transcurrido es cero.
Galileo estudió la caída de los cuerpos con experimentos de planos
inclinados y mostró que la naturaleza podía describirse con
matemáticas~\cite{galileo1623}. Newton, que nació el año en que murió
Galileo, reconoció la deuda con quienes lo precedieron:

% verificado: https://digitallibrary.hsp.org/index.php/detail/objects/9792 — original: «If I have seen further it is by standing on ye sholders of Giants»
\begin{cita}{Isaac Newton, carta a Robert Hooke (5 feb. 1675/76)~\cite{newton1676}}
Si he visto más lejos es porque estoy de pie sobre hombros de gigantes.
\end{cita}

\subsubsection*{Newton, Leibniz y dos notaciones}

Newton llamó a su método \emph{fluxiones} y escribía $\dot{y}$ para la
rapidez de cambio. Leibniz escribió $\dv{y}{x}$, pensando en un cociente
de cambios muy pequeños. Hoy usamos las dos, y también la notación $f'(x)$
de Lagrange. La obra principal de Newton, los \emph{Principia} (1687),
llegó a buena parte de Europa gracias a \textbf{Émilie du Châtelet}, que la
tradujo al francés y le añadió un comentario propio~\cite{chatelet1759}.

\subsubsection*{La derivada como lenguaje de la ciencia y la tecnología}

En 1843, \textbf{Ada Lovelace} describió cómo una máquina podría seguir
instrucciones para calcular, casi un siglo antes de las computadoras:

% verificado: https://psychclassics.yorku.ca/Lovelace/lovelace.htm — Nota A: «the Analytical Engine weaves algebraical patterns just as the Jacquard-loom weaves flowers and leaves»
\begin{cita}{Ada Lovelace, \emph{Notas} a la memoria de Menabrea (1843), Nota A~\cite{lovelace1843}}
La Máquina Analítica teje patrones algebraicos, así como el telar de
Jacquard teje flores y hojas.
\end{cita}

En el siglo XX, \textbf{Emmy Noether} demostró que cada simetría de las
leyes de la física corresponde a una cantidad que se conserva, un
resultado construido con derivadas. Al morir ella, Albert Einstein
escribió:

% verificado: https://mathshistory.st-andrews.ac.uk/Obituaries/Noether_Emmy_Einstein/ — «Pure mathematics is, in its way, the poetry of logical ideas»
\begin{cita}{Albert Einstein, carta a \emph{The New York Times} (1935)~\cite{einstein1935}}
La matemática pura es, a su manera, la poesía de las ideas lógicas.
\end{cita}

Y en 1962, \textbf{Katherine Johnson} verificó a mano, en la NASA, los
cálculos de la trayectoria orbital de John Glenn: posiciones, velocidades
y sus razones de cambio~\cite{shetterly2016}. La formalización moderna del
concepto que usaremos se encuentra en cualquier texto de cálculo
universitario~\cite{stewart2018}.

% ---------------------------------------------------------------------
% 4. Temas
% ---------------------------------------------------------------------
\tema{La pendiente: medir el cambio}\label{tema:pendiente}

\begin{hilo}
Galileo deja caer esferas por un plano inclinado y mide cuánto recorren
en cada intervalo de tiempo. Descubre algo sorprendente: la distancia no
crece de manera uniforme, sino que se acelera. ¿Cómo medir qué tan
rápido va la esfera?
\end{hilo}

La \textbf{razón de cambio promedio} de una función $f$ entre $x=a$ y
$x=b$ es la pendiente de la recta que une los dos puntos (recta secante):
\[
  m = \frac{\Delta y}{\Delta x} = \frac{f(b)-f(a)}{b-a}.
\]
Si $f(t)$ es una posición y $t$ el tiempo, esa pendiente es la
\emph{velocidad media}.

\begin{ejemplo}[Velocidad media]
Un objeto en caída libre recorre aproximadamente $d(t)=5t^2$ metros en
$t$ segundos. Entre $t=1$ y $t=3$:
\[
  v_{\text{media}} = \frac{d(3)-d(1)}{3-1} = \frac{45-5}{2} = \qty{20}{m/s}.
\]
\end{ejemplo}

\begin{aplicacion}[Ingeniería civil]
La pendiente de una rampa, una carretera o un tejado es una razón de
cambio: cuántos metros sube por cada metro que avanza. Los ingenieros la
calculan para que una silla de ruedas pueda subir una rampa o para que el
agua lluvia corra por un techo.
\end{aplicacion}

\begin{preguntas}
\pregunta Calcula la velocidad media del objeto anterior entre $t=2$ y
$t=4$. ¿Es mayor o menor que entre $t=1$ y $t=3$? ¿Por qué?
\espacio{2.5cm}
\pregunta Una rampa sube \qty{0,6}{m} en una distancia horizontal de
\qty{7,5}{m}. Calcula su pendiente y exprésala en porcentaje.
\espacio{2.5cm}
\end{preguntas}

\begin{resumen}
\begin{itemize}
  \item La razón de cambio promedio es $\dfrac{\Delta y}{\Delta x}$: la
        pendiente de la recta secante.
  \item Si la función es una posición, esa razón es la velocidad media.
  \item Una razón de cambio siempre tiene unidades: \unit{m/s}, \unit{\degreeCelsius\per\minute}, \ldots
\end{itemize}
\end{resumen}

% ---------------------------------------------------------------------
\tema{Del límite a la recta tangente}\label{tema:tangente}

\begin{hilo}
Newton, encerrado en su casa durante la peste de 1665, se hace la
pregunta clave: ¿qué pasa con la velocidad media cuando el intervalo de
tiempo se hace cada vez más pequeño, casi cero?
\end{hilo}

Si acercamos el punto $b$ al punto $a$, la recta secante se convierte en
la \textbf{recta tangente}. Escribiendo $b=a+h$:
\[
  f'(a) = \lim_{h\to 0} \frac{f(a+h)-f(a)}{h}.
\]
Este número es la \textbf{derivada} de $f$ en $a$: la pendiente de la
tangente y la razón de cambio \emph{instantánea}.

\begin{ejemplo}[Tangente a $f(x)=x^2$ en $x=1$]
\[
  f'(1) = \lim_{h\to0}\frac{(1+h)^2-1}{h}
        = \lim_{h\to0}\frac{2h+h^2}{h}
        = \lim_{h\to0}(2+h) = 2.
\]
La recta tangente en $(1,1)$ es $y = 2x-1$.

\begin{center}
\begin{tikzpicture}
\begin{axis}[mmc, xmin=-0.5, xmax=2.6, ymin=-1.2, ymax=5,
             xtick={1,2}, ytick={1,4},
             width=0.5\linewidth, height=0.42\linewidth,
             legend pos=outer north east,
             legend style={draw=none, font=\footnotesize}]
  \addplot+[domain=-0.3:2.2] {x^2};
  \addplot+[domain=0:2.5] {3*x-2};
  \addplot+[domain=0:2.5] {2*x-1};
  \addplot[only marks, mark=*, mark size=1.8pt] coordinates {(1,1) (2,4)};
  \legend{$y=x^2$, secante, tangente}
\end{axis}
\end{tikzpicture}
\end{center}
\end{ejemplo}

\begin{aplicacion}[Tecnología del celular]
El GPS de un celular no mide la velocidad directamente: registra tu
posición en intervalos muy cortos y calcula el cociente
$\Delta\text{posición}/\Delta t$. Mientras más corto el intervalo, más se
acerca a la velocidad instantánea: exactamente la idea del límite.
\end{aplicacion}

\begin{preguntas}
\pregunta Usa la definición con límite para hallar $f'(2)$ si $f(x)=x^2$.
Escribe la ecuación de la recta tangente en $(2,4)$.
\espacio{3cm}
\pregunta Completa la tabla para $f(x)=x^2$, $a=1$, y explica hacia qué
número se acerca el cociente:
\begin{center}
\begin{tabular}{c|cccc}
  $h$ & $0{,}1$ & $0{,}01$ & $0{,}001$ & $0{,}0001$\\\hline
  $\dfrac{f(1+h)-f(1)}{h}$ & & & & \\[2ex]
\end{tabular}
\end{center}
\end{preguntas}

\begin{resumen}
\begin{itemize}
  \item La derivada es el límite de las razones de cambio promedio:
        $f'(a) = \lim\limits_{h\to 0} \frac{f(a+h)-f(a)}{h}$.
  \item Geométricamente, $f'(a)$ es la pendiente de la recta tangente en
        $(a, f(a))$.
  \item Físicamente, si $f$ es una posición, $f'(a)$ es la velocidad instantánea.
\end{itemize}
\end{resumen}

% ---------------------------------------------------------------------
\tema{La derivada y sus reglas}\label{tema:reglas}

\begin{hilo}
Calcular límites para cada función es lento. Leibniz busca atajos: reglas
que den la derivada directamente. Tres siglos después, Katherine Johnson
usará esas mismas reglas para calcular, a mano, el regreso de un
astronauta a la Tierra.
\end{hilo}

\begin{teorema}[Reglas básicas de derivación]
Para constantes $c$ y $n$:
\[
  \dv{}{x}\,c = 0, \qquad
  \dv{}{x}\,x^n = n\,x^{n-1}, \qquad
  \dv{}{x}\bigl[c\,f(x)\bigr] = c\,f'(x), \qquad
  \dv{}{x}\bigl[f(x)\pm g(x)\bigr] = f'(x)\pm g'(x).
\]
\end{teorema}

\begin{ejemplo}[Un brazo robótico]
La posición (en \unit{cm}) de la pinza de un brazo robótico es
$s(t) = t^3 - 6t^2 + 9t$, con $t$ en segundos. Su velocidad es
\[
  v(t) = s'(t) = 3t^2 - 12t + 9 = 3(t-1)(t-3),
\]
así que la pinza se detiene ($v=0$) en $t=\qty{1}{s}$ y en $t=\qty{3}{s}$:
son los momentos en que cambia de dirección.
\end{ejemplo}

\begin{aplicacion}[Robótica]
Los robots corrigen su movimiento con un controlador que mira el error
(cuánto les falta para llegar), cuánto error han acumulado y la
\emph{derivada del error}: qué tan rápido se están acercando. Sin la
derivada, el robot frenaría tarde y se pasaría del objetivo.
\end{aplicacion}

\begin{preguntas}
\pregunta Deriva:
\begin{opciones*}(3)
  \item $f(x)=4x^5-3x^2+x-8$
  \item $g(t)=\dfrac{t^4}{2}-7t$
  \item $h(x)=\sqrt{x}$
\end{opciones*}
\espacio{3cm}
\pregunta Con la función $s(t)$ del brazo robótico, calcula la velocidad
en $t=\qty{2}{s}$. ¿Qué significa que sea negativa?
\espacio{3cm}
\end{preguntas}

\begin{resumen}
\begin{itemize}
  \item Regla de la potencia: $(x^n)' = n\,x^{n-1}$.
  \item La derivada de una suma es la suma de las derivadas; las
        constantes multiplicativas se conservan y las constantes solas se anulan.
  \item Donde $v(t)=s'(t)=0$, el objeto se detiene y puede cambiar de dirección.
\end{itemize}
\end{resumen}

% ---------------------------------------------------------------------
% 5. Cierre del hilo conductor
% ---------------------------------------------------------------------
\seccion{Cierre: de los hombros de gigantes a la Luna}

Newton y Leibniz discutieron durante años sobre quién había inventado
primero el cálculo. Hoy sabemos que ambos llegaron a él de forma
independiente, y que su invento fue más grande que su pelea: gracias a
du Châtelet se difundió, gracias a Noether se convirtió en lenguaje de la
física, y gracias a Katherine Johnson un ser humano pudo volver a salvo
del espacio. Ahora tú también estás sobre esos hombros.

% ---------------------------------------------------------------------
% 6. Evaluación (secciones institucionales)
% ---------------------------------------------------------------------
\matrizevaluacion
  {Interpreta la derivada como razón de cambio y como pendiente de la tangente.}
  {Calcula derivadas de polinomios y las usa para analizar el movimiento.}
  {Participa con curiosidad en la reconstrucción histórica del cálculo.}
  {Reconoce el aporte de mujeres y hombres de distintas épocas a la ciencia.}

\autoevaluacion

\seguimientodocente

% ---------------------------------------------------------------------
% 7. Referencias (toda cita del documento debe estar aquí)
% ---------------------------------------------------------------------
\begin{referencias}
% verificado: https://www.e-rara.ch/doi/10.3931/e-rara-14635 — pie de imprenta: Desaint & Saillant; Lambert, 2 vols., 1756–1759
\bibitem{chatelet1759} Du Châtelet, É. (1759). \emph{Principes mathématiques
  de la philosophie naturelle} (2 vols., 1756--1759). París: Desaint \& Saillant; Lambert.
% verificado: https://www.mathwomen.agnesscott.org/women/EinsteinNYTLetter.pdf — NYT, 4 may. 1935 (carta fechada el 1 de mayo)
\bibitem{einstein1935} Einstein, A. (1935, 4 de mayo). The late Emmy
  Noether [Carta al editor]. \emph{The New York Times}.
% verificado: https://portalegalileo.museogalileo.it/egjr.asp?c=36261 — Roma: Giacomo Mascardi, 1623
\bibitem{galileo1623} Galilei, G. (1623). \emph{Il Saggiatore}. Roma:
  Giacomo Mascardi.
% verificado: https://www.scss.tcd.ie/Brian.Coghlan/repository/J_Byrne/A_Lovelace/J_Fuegi_&_J_Francis_2003.pdf — Scientific Memoirs 3 (1843), pp. 666–731: traducción de Menabrea con las notas de Lovelace
\bibitem{lovelace1843} Lovelace, A. A. (1843). Sketch of the Analytical Engine
  invented by Charles Babbage, by L. F. Menabrea, with notes by the translator.
  \emph{Scientific Memoirs}, 3, 666--731.
% verificado: https://gblumen.mineducacion.gov.co/cgi-bin/koha/opac-detail.pl?biblionumber=4785 — MEN 2016, ISBN 978-958-691-925-8
\bibitem{mendba} Ministerio de Educación Nacional (2016). \emph{Derechos
  Básicos de Aprendizaje V.2: Matemáticas}. Bogotá: MEN.
% verificado: https://digitallibrary.hsp.org/index.php/detail/objects/9792 — original en la Historical Society of Pennsylvania, colección Simon Gratz
\bibitem{newton1676} Newton, I. (1675/76, 5 de febrero, calendario juliano). Carta a
  Robert Hooke. Historical Society of Pennsylvania, colección Simon Gratz.
% verificado: https://cmc.marmot.org/Record/.b62339163 — William Morrow, 2016
\bibitem{shetterly2016} Shetterly, M. L. (2016). \emph{Hidden Figures}.
  Nueva York: William Morrow.
% verificado: https://www.cengage.com/c/ebook-c%C3%A1lculo-de-una-variable-trascendentes-tempranas-8e-stewart/9780357087244/ — Cengage, 8.ª ed.
\bibitem{stewart2018} Stewart, J. (2018). \emph{Cálculo de una variable:
  trascendentes tempranas} (8.ª ed.). México: Cengage Learning.
\end{referencias}

\end{document}
